\section {Background and overview}

{\tt LatMRG} is a software system implemented as a library of modules
written in the Modula-2 language.
It provides different tools for studying the structure of lattices 
in the real space and
for examining the theoretical properties of random number
generators based on linear recurrences in modular arithmetic.
It offers facilities for checking if a generator has maximal period or not,
for examining its lattice structure 
(e.g., applying lattice and spectral tests),
and for performing computer searches for ``good'' generators
according to different quality criteria.
The software can also be used for related applications, such as 
searching and evaluating lattice rules for quasi-Monte Carlo integration.

In this section, we give a quick recall of some definitions and
notation, as well as a short outline of what the package does.
For more details on the underlying theory and algorithms,
see \cite{rLEC97c} and other references given there.
We classify the modules of {\tt LatMRG} in three groups:
(a) low-level, (b) intermediate-level, and (c) high-level.
Higher-level modules import facilities from the lower-level ones.

The high-level modules (c) are programs in executable form
which read their data in files or work interactively with the user.
They can either analyze a given generator or seek ``good'' generators
according to different criteria.
Examples of data files and results are given in Section 2.
Appendix A gives specifications of the data file formats 
and of what the programs do.

The intermediate-level modules (b) provide data types and procedures
to construct lattice bases for
different classes of generators (simple or combined MRGs, lacunary
indices, etc.), manipulate such bases, find a shortest vector in a
lattice, reduce a basis in the sense of Minkowski, and so on.
These tools are used by the upper-level modules (c), but can also be
used directly to make programs different than those already provided
at level (c), offering thus more flexibility.
The lower-level modules (a) implement basic operations on scalars, vectors,
matrices, polynomials, and so on.
They allow different possible representations for these objects,
depending, for example, on the size of the modulus $m$ and 
the precision we want, as explained in Section 1.6.
These lower-level tools are used by the modules of levels (b) and (c).
The intermediate and low-level modules are discussed a little further
in Section 3, and their specifications are given in appendices B and C.


%%%%%%%%%%%%%%%%%%%%%%%%%%%%%%%%%%%%%%%%%%%%%%%%%%%%%%%%
\subsection {Lattices in the real space}

The {\em lattices\/} considered here are discrete subspaces of the 
real space $\R^t$, which can be expressed as 
\eq
 L_t = \left\{\bv = \sum_{j=1}^t z_j\bv_j\mid \mbox{ each } z_j\in\Z\right\},
                                                       \label {eq:lattice}
\endeq
where $t$ is a positive integer, and
$\bv_1,\dots,\bv_t$ are linearly independent vectors in $\R^t$
which form a {\em basis\/} of the lattice.
A comprehensive treatment of such lattices can be found in \cite{mCON99a}.
The matrix $\bV$, whose $i$th line is $\bv_i$, is the
corresponding {\em generator matrix\/} of $L_t$.
A lattice $L_t$ shifted by a constant vector $\bv_0\not\in L_t$,
i.e., a point set of the form $L'_t = \{\bv+\bv_0 : \bv \in L_t\}$,
is called a {\em grid\/}, or a {\em shifted lattice\/}.
The lattices considered in this guide always contain, or are contained in,
the integer lattice $\Z^t$, i.e., $\Z^t\subseteq L_t$ or $L_t\subseteq \Z^t$.

The {\em dual lattice} of $L_t$ is defined as
$L_t^* = \{\bh \in \R^t : \bh\cdot\bv \in\Z$ for all $\bv\in L_t\}$.
The {\em dual\/} of a given basis $\bv_1,\dots,\bv_t$ is the set of
vectors $\bw_1,\dots,\bw_t$ in $\R^t$
such that $\bv_i\cdot\bw_j = \delta_{ij}$, where $\delta_{ij}=1$
if $i=j$, and $\delta_{ij}=0$ otherwise.
It is a basis of the dual lattice.
These $\bw_j$'s are the columns of the matrix $\bV^{-1}$,
the inverse of the matrix $\bV$.
If $m$ is any positive real number, a basis $\{\bw_1,\dots,\bw_t\}$
satisfying $\bv'_i \bw_j = \delta_{ij} m$ for all $i$, $j$ is called
the $m$-dual of the basis $\{\bv_1,\dots,\bv_t\}$. 
The lattice generated by this $m$-dual basis is the $m$-dual to $L$. 
This extension of the usual notion of dual basis and dual lattice
will allow us, by a suitable choice of $m$ (in our context it will be 
the modulus in (\ref{recurs})), to deal uniquely with integer coordinate 
vectors,  which can be represented exactly on a computer.

The determinant of the matrix $\bV$ is equal to the volume of the
fundamental parallelepiped
$\Lambda = \{\bv = \lambda_1\bv_1 + \cdots + \lambda_t\bv_t :
0\le \lambda_i\le 1$ for $1\le i\le t\}$, and is also the inverse
of the average number of points per unit of volume, independently of 
the choice of basis. It is called the determinant of $L_t$.
The quantity $1/\det(L_t) = 1/\det(\bV) = \det(\bV^{-1})$
is called the {\em density\/} of $L_t$.
When $L_t$ contains $\Z^t$, the density is an integer equal to the 
cardinality of the point set $L_t \cap [0,1)^t$.

For a given lattice $L_t$ and a subset of coordinates
$I = \{i_1,\dots,i_d\} \subseteq \{1,\dots,t\}$, denote by $L_t(I)$
the projection of $L_t$ over the $d$-dimensional subspace determined
by the coordinates in $I$.
This projection is also a lattice, whose density divides that of $L_t$.
There are exactly $\det(L_t(I))/\det(L_t)$ points of $L_t$ that are
projected onto each point of $L_t(I)$.
In group theory language, $L_t(I)$ corresponds to a coset of $L_t$.

%%%%%%%%%%%%%%%%%%%%%%%%%%%%%%%%%%%%%%%%%%%%%%%%%%%%%%%%%
\subsection {Multiple recursive generators}

Consider the linear recurrence
\eq
   x_n = (a_1 x_{n-1} + \cdots + a_k x_{n-k}) \mod m.       \eqlabel{recurs}
\endeq
where $m$ and $k$ are positive integers and each $a_i$ belongs to the
set (or ring) $\Z_m = \{0,1,\dots,m-1\}$.
For $n\ge 0$, $s_n = (x_{n},\dots,x_{n+k-1}) \in \Z_m^k$ is the {\em state}
at step $n$.  The initial state $s_0$ is called the {\em seed}.
One can take $u_n = x_n/m \in [0,1)$ as the {\em output} at step $n$.
This kind of generator is called {\em multiple recursive} (MRG).
When $k=1$, it gives the well-known multiplicative linear congruential
generator (MLCG).
MLCGs in {\em matrix form\/} can also be expressed as
many copies of the same MRG running in parallel.
For more details, see \cite{rGRO88a,rLEC90a,rLEC94a,rNIE86a}.

The maximal possible period for the $s_n$'s is the
cardinality of $\Z_m^t$ minus 1, i.e. $\rho = m^k-1$.
It is attained if and only if $m$ is prime and
the characteristic polynomial of (\ref{recurs}),
\eq
  P(z) = \left(z^k - \sum_{i=1}^k a_i z^{k-i}\right) \mod m,  \eqlabel{charac}
\endeq
is a primitive polynomial modulo $m$.
Knuth \cite {rKNU97a} gives necessary and sufficient
conditions for that, which are implemented in our package.
If $k=1$ and $m=p^e$, with $e > 1$, then
the maximal possible period is $2^{e-2}$ for $p=2$
and $(p-1)p^{e-1}$ for $p > 2$ \cite{rKNU97a,rLEC90a}.

Instead of taking $u_n = x_n/m$ for the output, one can take a more
general linear combination of the components of the state vector, say
\begin {eqnarray}
 y_n &=& (b_1 x_{n} + \cdots + b_k x_{n+k-1}) \mod m,    \eqlabel {yn} \\
 u_n &=& y_n /m.                                         \eqlabel {output2}
\end {eqnarray}
For any integer $t\ge 1$, one has
\eq
    \pmatrix{ y_n \cr y_{n+1} \cr \vdots\cr y_{n+t-1}}
  = \pmatrix{ \bb' \cr \bb'\bA \cr \vdots\cr \bb'\bA^{t-1}}
    \pmatrix{ x_{n} \cr x_{n+1} \cr \vdots\cr x_{n+k-1}} \; \mod m
  \eqdef \bB_t
    \pmatrix{ x_{n} \cr x_{n+1} \cr \vdots\cr x_{n+k-1}} \; \mod m,
                                                   \eqlabel {eq:tpoly}
\endeq
where $\bb' = (b_1,\dots,b_k)$ and
$$
  A = \pmatrix{
     0   &1      &\ldots &0 \cr
  \vdots &\vdots &\ddots &\vdots\cr
     0   &0      &\ldots &1 \cr
     a_k &a_{k-1}&\ldots &a_1\cr}                 \eqlabel{eq:companion}
$$
is the {\em companion matrix\/} of the characteristic polynomial $P(z)$.
In particular, by taking $t=k$, one sees that the vector 
$(y_n,\dots,y_{n+k-1})$ takes all possible values 
in $\Z_m^k$, when $(x_n,\dots,x_{n+k-1})$ does so, if and only if
the matrix $\bB_k$ has full rank.
The matrix $\bB_t$ can be constructed easily as follows.
Put $(x_n,\dots,x_{n+k-1})' = \be_j$, the $j$th vector of the canonical
basis, with $x_{n+i-1} = \delta_{ij}$, and compute the corresponding 
column vector $(y_n,\dots,y_{n+t-1})'$ via (\ref{recurs}) and (\ref{yn}).
This vector is the $j$th column of the matrix $\bB_t$.

%%%%%%%%%%%%%%%%%%%%%%%%%%%%%%%%%%%%%%%%%%%%%%%%%%%%%%%%%%%%%%%%
\subsection {Lattice structure and spectral test}

Let $\Psi_t$ be the multiset of all the $t$-dimensional vectors of 
successive output values of an MRG, from all possible seeds in $\Z_m^k$, i.e.,
$$
 \Psi_t = \{\bu_{0,t} = (x_0/m,\dots,x_{t-1}/m) : (x_0,\dots,x_{k-1})\in \Z_m^k\}.
$$
For $t\le k$, this set is just $\Z_m^t$ with each element repeated
$m^{k-t}$ times.
For $t > k$, the first $k$ components of a vector 
$\bu_{0,t}\in\Psi_t$ are arbitrary elements of $\Z_m/m$,
but once they are fixed, the remaining $t-k$ components 
are determined uniquely by the linear recurrence (\ref{recurs}).
The last $t-k$ components are thus linear combinations modulo $1$,
with integer coefficients, of the first $k$ components.

For $1\le i\le k$, let $\bv_i = (v_{i,1},\dots,v_{i,t})$ be the 
$t$-dimensional vector with components $v_{i,j}=\delta_{ij}/m$
for $i\le k$, and $v_{i,j} = (a_1 v_{i,j-1} +\cdots + a_k v_{i,j-k})\mod 1$
for $j > k$.
For $k+1\le i\le t$, let $\bv_i = \be_i$, the $i$th unit vector in $t$
dimensions.  These vectors are a basis of a lattice $L_t$ that contains
$\Z^t$, with unit cell volume of $\max(m^{-t}, m^{-k})$, such that
$L_t \cap [0,1)^t = \Psi_t$.  In fact,
$\L_t = \Psi_t + \Z^t = \{\bv = \tilde\bv + \bz : \tilde\bv\in\Psi_t$
and $\bz\in\Z^t\}$.
The vectors $\bw_i = (w_{i,1},\dots,w_{i,t})$, $1\le i\le t$, where 
$$
  w_{i,j} = \cases {m  & for $j=i\le k$;\cr
                    0  & for $j\not=i\le k$;\cr
                    v_{j,i} & for $i > k\ge j$;\cr
                    1       & for $j=i > k$;\cr
                    0       & for $k < j\not=i >k$,\cr }
$$
are linearly independent and satisfy $\bv_i\cdot\bw_j = \delta_{ij}$.   
They form the dual basis to $\{\bv_1,\dots,\bv_t\}$.
The vectors $\bv_i$ and $\bw_j$ are the lines of the matrices:
$$
  \bV_t
  = (\bv_1 \bv_2 \cdots \bv_t)'
  = \pmatrix{
     1/m &0      &\ldots &0   & v_{1,k+1} & \ldots & v_{1,t} \cr
     0   & 1/m   &\ldots &0   & v_{2,k+1} & \ldots & v_{2,t} \cr
  \vdots &\vdots &\ddots &\vdots &\vdots  &        & \vdots  \cr
     0   & 0     &\ldots &1/m & v_{k,k+1} & \ldots & v_{k,t} \cr
     0   & 0     &\ldots  & 0  & 1      & \ldots & 0       \cr
  \vdots &\vdots &\ddots &\vdots &\vdots  & \ddots & \vdots  \cr
     0   & 0     &\ldots  & 0  & 0        & \ldots & 1     \cr }
$$
and 
$$
  \bW_t
  = (\bw_1 \bw_2 \cdots \bw_t)'
  = \pmatrix{
     m   & 0     &\ldots &0   & 0 & \ldots & 0 \cr
     0   & m     &\ldots &0   & 0 & \ldots & 0 \cr
  \vdots &\vdots &\ddots &\vdots &\vdots & & \vdots  \cr
     0   & 0     &\ldots &m   & 0 & \ldots & 0 \cr
 -v_{1,k+1}& -v_{2,k+1}&\ldots& -v_{k,k+1} & 1 & \ldots & 0  \cr
  \vdots &\vdots &\ddots &\vdots &\vdots  & \ddots & \vdots  \cr
 -v_{1,t}& -v_{2,t}&\ldots& -v_{k,t} & 0  & \ldots & 1  \cr },
$$
and one has $\bW'_t \bV_t = I$.
In the package {\tt LatMRG}, we store the vectors $m\bv_i$
instead of $\bv_i$ in the computer, for the components of the 
former vectors are integer-valued and can thus be represented 
exactly in the computer.


For the more general case of (\ref{yn}) and (\ref{output2}),
replace the first $k$ lines of $\bV_t$ by $\bB'_t$.
If $\bB_k$ is invertible, then $\bB_t$ has rank $k$ and the lines 
$k+1$ to $t$ of $\bV$ complete the lattice basis as before.
Otherwise, remove the lines in $\bB'_t$ which are linearly dependent 
of others, to obtain a matrix of full rank $k'<k$, and replace them by
$k-k'$ vectors of the canonical basis of $\R^t$, divided by $m$, chosen 
in a way that the first $k$ lines and $k$ columns of $\bV_t$ form an invertible
matrix.
In both cases, the dual basis is obtained by inverting the matrix $\bV'_t$.
The module {\tt MRG} does that.
For LCGs in matrix form, bases for $L_t$ and 
its dual can be constructed as explained in \cite{rGRO88a,rLEC97c}.  

If one adds a constant $b$ on the right-hand-side of (\ref{recurs}),
before applying the modulo operation, then the vectors of successive values
will all belong to $L'_t$, where $L'_t = L_t + \bv_{0,t}$ is a shift of $L_t$
by some constant $v_{0,t}\in\Z_m^t$, i.e., a {\em grid\/}.
Since $L'_t$ and $L_t$ have the same structural properties, we simply
ignore the presence of such a constant $b$ in {\tt LatMRG}, and 
consider only homogeneous recurrences.

When $m$ is prime and the MRG has full period $m^k-1$, then $\Psi_t$ is the
set of all $t$-tuples produced by the generator over its main cycle,
plus the zero vector.
Otherwise, the set of
$t$-dimensional vectors produced over any given (sub)cycle (plus the
zero vector and plus $m\Z^t$) is a strict subset of $L_t$ which in
general does not form a lattice.
Then, {\tt LatMRG} can analyze the set of all $t$-tuples produced over
the {\em union\/} of all subcycles.
In some cases, however, the vectors of
successive values over one subcycle generate a strict sublattice
of $L_t$, whose intersection with $[0,1)^t$ contains only a fraction of
the points of $\Psi_t$.
This is what happens in particular when $k=1$, $m$ is a power of a
prime $p$, and $x_0$ is prime to $p$.
The package can take care of the latter case by constructing a
basis for the appropriate sublattice.
%% For instance, that sublattice contains one-fourth of the points when $m$
%% is a power of 2, $a_1 \equiv 5 (\mod 8)$, and $x_0$ is odd \cite{rKNU81a}.

%%%%%%%%%%%%%%%%%%%%%%%%%%%%%%%%%%%%%%%%%%%%%
\subsection {Lacunary indices}

Instead of forming vectors with successive values like in the above
definition of $\Psi_t$,
one can form vectors with values that are some distance apart in the
sequence (so-called ``leapfrog'' values).
Let $I=\{i_1 , i_2 , \cdots , i_t\}$ be a set of fixed integers.
Define
\eq
 \psi_t(I) = \{(u_{i_1},\dots,u_{i_t}) \mid
             (x_0,\dots,x_{n+k-1}) \in \Z_m^k\}                  \eqlabel {psi2}
\endeq
and let $L_t(I) = \psi_t(I) + \Z^t$.
If we assume that $0\le i_1 < i_2 < \cdots < i_t$, 
this $L_t(I)$ is the projection of the lattice $L_{i_t+1}$ over the $t$-dimensional
subspace determined by the coordinates that belong to $I$.
Using the module {\tt MRG}, one can build a basis for $L_t$ and its
dual in this more general case, and then perform lattice analysis as usual.
Further details and examples are given in \cite{rLEC97c}.
For $(i_1,\dots,i_t) = (0,\dots,t-1)$, one has $L_t(I) = L_t$.
%  Distances between hyperplanes, Beyer quotients, etc.,
%  can then be computed as usual.

To construct the basis in this case, one must compute the vector
$(u_{i_1},\dots,u_{i_t})$ obtained when the seed $(x_0,\dots,x_{k-1})=\be_i$,
for each vector $\be_i$ of the canonical basis.
The linear transformation from the state $(x_0,\dots,x_{k-1})$ to the 
vector $(u_{i_1},\dots,u_{i_t})$ is one-to-one for each $t\ge k$ if and only if
the transformation applied to the $k$ vectors of the canonical basis gives
$k$ linearly independent vectors for $t=k$.
For $t < k$, the transformation is onto (surjective) if and only if the 
transformation gives $t$ linearly independent vectors; that is, the corresponding
matrix has full rank $t$.

%%%%%%%%%%%%%%%%%%%%%%%%%%%%%%%%%%%%%%%%%%%%%
\subsection {Figures of Merit}

Figures of merit measure the quality of lattices.
Here, good quality means that the points cover the space very evenly,
i.e., are very uniformly distributed.
There are many ways of measuring this uniformity, which give rise
to several different figures of merit.

The lattice structure also means that all points of $L_t$ lie in a family
of equidistant parallel hyperplanes.  Among all such families of
hyperplanes that cover all the points, choose the one for which the successive
hyperplanes are farthest apart.
The distance between these successive hyperplanes is in fact equal to 
$1/\ell_t$ where $\ell_t$ is the Euclidean 
length of the shortest nonzero vector in the {\em dual} lattice $L_t^*$.
So for a given density of points, we want $\ell_t$ to be as 
large as possible.  Computing this $\ell_t$ for an MRG and comparing
with the best possible value, given $t$, $m$, and $k$, is known
as the {\em spectral test\/} in the literature on RNGs 
\cite{rKNU97a,sFIS96a}.

We can view the lattice as a way of packing the space by spheres of
radius $\ell_t/2$, with one sphere centered at each lattice point.
In the dual lattice, this gives $1/n = m^{-k}$ spheres per unit of volume.
If we rescale so that the radius of each sphere is 1,
we obtain $\delta_t = (\ell_t/2)^t/n$ spheres per unit of volume.
This number $\delta_t$ is called the {\em center density\/} of the lattice.
For a given value of $n$, an upper bound on $\ell_t$ can be obtained 
in terms of an upper bound on $\delta_t$ (one has 
$\ell_t = 2 (n\delta_t)^{1/t}$), and vice-versa.
Let $\delta_t^*$ be the largest possible value of $\delta_t$ for a
lattice (i.e., the densest packing by non-overlapping spheres arranged
in a lattice).  The quantity $\gamma_t = 2(\delta_t^*)^{2/t}$
is called the {\em Hermite constant\/} for dimension $t$
\cite{mCON99a,mGRU87a}.  It gives the upper bound 
$\ell_t^2\le (\ell_t^*(n))^2 = 2(n\delta_t^*)^{2/t} = \gamma_t n^{2/t}$
for a lattice of density $1/n$.
Knowing the Hermite constants, or good approximations of them, is useful
because it allows us to normalize $\ell_t$ to a value between 0 and 1 by
taking $\ell_t/\ell_t^*(m^k)$.  This is convenient for comparing values for
different values of $t$ and $m^k$.
Good values are close to 1 and bad values are close to 0.

The Hermite constants are known exactly only for $t\le 8$, in which 
case the densest lattice packings are
attained by the {\em laminated\/} lattices \cite{mCON99a}.
Conway and Sloane \cite[Table 1.2]{mCON99a} give the 
values of $\delta_t^*$ for $t\le 8$, and provide lower and upper
bounds on $\delta_t^*$ for other values of $t$.
The largest value of $\ell_t^2/n^{2t}$ obtained so far for concrete lattice 
constructions is a lower bound on $\gamma_t$, which we denote by
$\gamma_t^{\rmB}$.  Such values are given in Table 1.2 of \cite{mCON99a}, 
page 15, in terms of $\delta^*$.
The laminated lattices, which give the lower bound 
$\ell_t^2/n^{2t} \ge \gamma_t^{\rmL} = 4 \lambda_t^{-1/t}$,
where the constants $\lambda_t$ are given in 
\cite[Table 6.1, page 158]{mCON99a} for $t\le 48$,
are the best constructions in dimensions 1 to 29, except for 
dimensions 10 to 13.
(One has $\gamma_t^{\rmL} = \gamma_t$ for $t\le 8$.)

Minkowski proved that there exists lattices with density satisfying
$\delta_t \ge \zeta(t) / (2^{t-1} V_t)$ where $\zeta(t) = \sum_{k=1}^\infty$
is the Riemann zeta function and $V_t = \pi^{n/2} / (n/2)!$ 
is the volume of a $t$-dimensional sphere of radius 1.
This bound provides a lower bound $\gamma_t^{\rm Z}$ on $\gamma_t$.

An upper bound on $\gamma_t$ is obtained via the bound of Rogers on the 
density of sphere packings \cite{mCON99a}.
This upper bound can be written as
$$
 \gamma_t^{\rmR} = 4 e^{2R(t)/t}
$$
where $R(t)$ can be found in Table~1.2 of \cite{mCON99a} for $t\le 24$,
and can be approximated with $O(1/t)$ error and approximately 4 decimal
digits of precision, for $t\ge 25$, by
\eq
 R(t) =   {1\over 2} n \,\log_2\left({n\over 4\pi e}\right)
        + {3\over 2} \log_2 (n)
        - \lg (e/\sqrt{\pi})  + {5.25\over n + 2.5}.
\endeq
Table~1 in \cite{rLEC99c} gives the ratio 
$(\gamma_t^{\rmL} / \gamma_t^{\rmR})^{1/2}$,
of the lower bound over the upper bound on $\ell_t$, for $1\le t\le 48$.
This ratio tends to decrease with $t$, but not monotonously.

Computing the shortest vector in terms of the Euclidean norm is convenient,
e.g., for computational reasons, but one can also use another norm instead.
For example, one can take the $\cL_p$-norm, defined by
$\Vert\bv\Vert_p = (|v_1|^p + \cdots + |v_t|^p)^{1/p}$ for $1\le p < \infty$
and $\Vert\bv\Vert_\infty = \max(|v_1|, \dots, |v_t|)$ for $p = \infty$.
The inverse of the length of the shortest vector is then the $\cL_p$-distance
between the successive hyperplanes for the family of hyperplanes that are 
farthest apart among those that cover $L_t$.
For $p=1$, the length $\ell_t = \Vert\bv\Vert_1$ of the shortest vector $\bv$
(or $\Vert\bh\Vert_1 - 1$ in some cases, see \cite{rKNU97a})
is the minimal number of hyperplanes that cover all the points of $\Psi_t$.
The following upper bound on $\ell_t$ in this case was established by 
Marsaglia \cite{rMAR68a} by 
applying the general convex body theorem of Minkowski:
$$ 
  \ell_t \le \ell_t^*(m^k) = (t! m^k)^{1/t} \eqdef \gamma_t^{\rm M} m^{k/t}.
$$
This upper bound can be used to normalize $\ell_t$ in this case.

%  --- Upper bound on $\ell_t$ in general: Minkowski.  ?

As a figure of merit, we take the worst-case value of $\ell_t/\ell_t^*(m^k)$
over certain values of $t$ and for selected projections on lower-dimensional
subspaces.  More specifically, let $\ell_I$ denote the length
of the shortest nonzero vector $\bv$ in $L_t^*(I)$,
and $\ell_t = \ell_{\{1,\dots,t\}}$ as before.
For arbitrary positive integers $t_1\ge \cdots\ge t_d \ge d$,
consider the worst-case figure of merit
\eq
 M_{t_1,\dots,t_d} = \min\left[ \min_{k+1\le t\le t_1} \ell_t/\ell_t^*(m^k),\,
    \min_{2\le s\le k} \min_{I\in S(s,t_s)} \ell_I/m, \,
    \min_{k+1\le s\le d} \min_{I\in S(s,t_s)} \ell_I/\ell_s^*(m^k) \right],
                                                            \label{eq:M}
\endeq
where $S(s,t_s) = \{I=\{i_1,\dots,i_s\} : 1 = i_1 < \cdots < i_s\le t_s\}$.
This figure of merit makes sure that the lattice is good in projections 
over $t$ successive dimensions for all $t\le t_1$, 
and over non-successive dimensions that are not too far apart.
Note that when $s\le k$, the smallest distance between hyperplanes that can
be achieved in $s$ dimensions for the MRG is $1/m$, so $\ell_s/m$ cannot
exceed 1, and it is equal to 1 if and only if the linear transformation
from the state $(x_0,\dots,x_{k-1})$ to the output vector 
$(u_{i_1},\dots,u_{i_s})$ is surjective (i.e., the corresponding matrix has
full rank).  For $s<k$, $m$ is typically much smaller than $\ell_s^*(m^k)$,
and this is the reason for separating the last two terms in (\ref{eq:M}).

The figure of merit $M_{t_1} = \min_{2\le s\le t_1} \ell_s/\ell_s^*(n)$
(with $d=1$) has been widely used for ranking and selecting
LCGs and MRGs \cite{sFIS96a,rLEC99b,rLEC99c}.
% \cite{vLEM99d} suggested using $d=2$ or 3 instead of $d=1$,
%  with $t_1 = \cdots = t_d$.
The quantity $M_{t_1,\dots,t_d}$ is a worst case over
$(t_1-d) + \sum_{s=2}^d {t_s-1\choose s-1}$ projections, and this
number increases quickly with $d$ unless the $t_s$ are very small.
For example, if $d=4$ and $t_s = t$ for each $s$,
there are 5019 projections for $t = 32$.
When too many projections are considered, there are inevitably some
that are bad, so the worst-case figure of merit is (practically)
always small, and can no longer distinguish
between good and mediocre behavior in the most important projections.
Moreover, the time to compute $M_{t_1,\dots,t_d}$ increases
with the number of projections.
We should therefore consider only the projections deemed important.
We suggest using the criterion (\ref{eq:M}) with $d$ equal to 4 or 5,
and $t_s$ decreasing with $s$.

Instead of considering the shortest nonzero vector in the dual lattice,
one can consider the shortest nonzero vector in the primal lattice $L_t$.
Its length represents the distance to the nearest other lattice point 
from any point of the lattice.  A small value means that many points
are placed on the same line, at some fixed distance apart.

%%%%%%%%%%%%%%
\subsubsection {Minkowski-reduced basis}

Another way of measuring the quality of a lattice is in terms of the
relative lengths of the smallest and largest vectors in a 
{\em reduced\/} basis.
A basis can be {\em reduced\/} in different senses.
One type of reduced basis considered by this package is a 
{\it Minkowski-reduced lattice basis\/} (MRLB)
(see \cite{rAFF85a,rAFF88a,rGRO88a} for more details).
Roughly, a MRLB is a basis for which the vectors are in some sense the 
most orthogonal.
The ratio of the sizes of the shortest and longest vectors of a MRLB
is called its {\it Beyer-quotient\/}.
In general, a given lattice may have several MRLBs, all with
the same length of the shortest vector, but perhaps with
different lengths of the longest vector, and thus different Beyer quotients.
We define $q_t(I)$ as the maximum of the Beyer quotients of all MRLBs
of $L_t(I)$, and denote $q_t(\{1,\dots,t\})$ by $q_t$.   
We prefer values of $q_t(I)$ close to 1.
Similar to (\ref{eq:M}), we define
\eq
 Q_{t_1,\dots,t_d} = \min\left[ \min_{k+1\le t\le t_1} q_t,\,
    \min_{2\le s\le d} \min_{I\in S(s,t_s)} q_t(I) \right].
                                                            \label{eq:Q}
\endeq

%%%%%%%%%%%%%%%%
\subsubsection {The $\cP_\alpha$ criterion}

The quantity $\cP_\alpha$ is a measure of non-uniformity 
(i.e., discrepancy from the uniform distribution, the smaller the better),
which has been widely used in the context of quasi-Monte Carlo integration
(see, e.g., \cite{vSLO94a}).
In the case where $\Psi_t = L_t \cap [0,1)^t$ where $L_t$ is a lattice
with dual $L_t^*$, one has
\eq
 \cP_\alpha(\Psi_t) =
   \sum_{\bzero\not=\bw\in L_t^*} \|\bw\|_{\pi}^{-\alpha},
                                                  \label{eq:palpha}
\endeq
where $\|\bw\|_{\pi}=\prod_{j=1}^t \max(1,|w_j|)$ for $\bw = (w_1,\dots,w_t)$.
For any positive integer $\alpha$, $\cP_{2\alpha}(\Psi)$ can be
written equivalently as
\begin{eqnarray}
 \cP_{2\alpha}(\Psi_t)
  = -1 + \frac{1}{n}\sum_{\bu\in \Psi_t}
    \prod_{j=1}^t \left[1 - \frac{(-4\pi^2)^{\alpha}}{(2\alpha) !}
    B_{2\alpha}(u_j)\right]                      \label{eq:palpha2}
\end{eqnarray}
where the $B_\alpha$ are the Bernoulli polynomials:
\begin {eqnarray*}
 B_0(x) &=& 1,\\
 B_1(x) &=& x -1/2,\\
 B_2(x) &=& x^3 - x^2 + 1/6,\\
 B_3(x) &=& x^3 - 3x^2/2 + x/2,\\
 B_4(x) &=& x^4 - 2x^2 + x - 1/30,
\end {eqnarray*}
and the other polynomials can be found via the identity
$$
  {t e^{xt} \over e^t-1} = \sum_{i=0}^\infty B_i(x) t^i/i!.
$$
Hickernell \cite{vHIC98c} introduced generalizations of $\cP_\alpha$,
incorporating weights and replacing the simple sum in (\ref{eq:palpha})
by a more general norm.  One version of this weighted $\cP_\alpha$,
where the weight associated to the projection over the coordinates
in a set $I$ has the {\em product-form\/}
$\beta_I = \beta_0\prod_{j\in I}\beta_j$, can be defined by
\[
  \cP_{2\alpha} (\Psi_t)
  = -\beta_0 + \frac{\beta_0}{n} \sum_{\bu\in \Psi_t}
    \prod_{j=1}^t \left[1 - \frac{(-4\pi^2\beta_j^2)^{\alpha}}{(2\alpha)!}
    B_{2\alpha}(u_j)\right]                       \label{eq:palpha3}
\]
when $\alpha$ is an integer.  The identity (\ref{eq:palpha2})
or (\ref{eq:palpha3}) gives an algorithm for computing
$\cP_{2\alpha}(\Psi_t)$ in time $O(nt)$ when $\alpha$ is an integer
and $\Psi_t$ is the intersection of a lattice with $[0,1)^t$.
Note that the $D_{\cF,\alpha,p}(P)$ of \cite{vHIC00b} corresponds to
$(\cP_{2\alpha}(\Psi_t))^{1/2}$ for $\Psi_t=P$, 
$p=2$, and $\beta_j=1$ for all $j$.
LatMRG provides tools for computing $\cP_\alpha$ with or without weights.

It is has been proved (e.g., \cite{vSLO94a}, Theorem 4.4, page 83)
that for any $t\ge 2$, $\alpha > 1$, and
prime number $m > e^{\alpha t/(\alpha-1)}$, there exists at least one
LCG with modulus $m$ such that 
\eq
  \cP_\alpha(\Psi_t) \le {[(e/t)(2\ln m + t)]^{\alpha t}\over m^\alpha}.
\endeq
The latter quantity can then be used to normalize $\cP_{\alpha}$.

Alternatively, Hickernell et al.\ \cite{vHIC01a}, section 4.1, 
suggest using the figure of merit $g_t$, where
\begin {eqnarray}
 g_t^2 &=& {n^2\over (3/2)^t-1} \left({t-1\over t-1+\log n}\right)^{t-1}
           \cP_{2}(\Psi_t) \nonumber \\
       &=& {n\over (3/2)^t-1} \left({t-1\over t-1+\log n}\right)^{t-1} 
       \left[-n + \sum_{\bu\in \Psi_t}
       \prod_{j=1}^t (1 + 3 B_{2\alpha}(u_j))\right]
                                                       \label{eq:palpha4}
\end {eqnarray}
is a normalized version of the inverse of (\ref{eq:palpha3}) 
with $\alpha=1$, $\beta_0=1$, and $\beta_j = \pi\sqrt{3/2}$ for $j\ge 1$.
This $g_t^2$ can be rewritten as
\begin {eqnarray}
 g_t^2 &=& \gamma_t^{\rmP}(n) 
       \left[-n + \sum_{\bu\in \Psi_t}
       \prod_{j=1}^t (1 + 3 B_{2\alpha}(u_j))\right]
                                                       \label{eq:palpha5}
\end {eqnarray}
where 
\eq
  \gamma_t^{\rmP}(n) 
  = {n\over (3/2)^t-1} \left({t-1\over t-1+\log n}\right)^{t-1} 
\endeq
is a constant that depends on $t$ and $n$.
As a figure of merit based on $\cP_2$, similar to (\ref{eq:M}), we define
\eq
 G_{t_1,\dots,t_d} = \min\left[ \min_{k+1\le t\le t_1} 1/g_t(\Psi_t),\,
%    \min_{2\le s\le k} \min_{I\in S(s,t_s)} 1/g_t(\Psi_t(I)), \,
    \min_{k+1\le s\le d} \min_{I\in S(s,t_s)} 1/g_t(\Psi_t(I)) \right].
                                                            \label{eq:G}
\endeq


%%%%%%%%%%%%%%%%%%%%%%%%%%%%%%%%%%%%%%%%%%%%
\subsection {Matrix multiple recursive generators}

MRGs in matrix form, which we denote MMRGs, have been introduced
and studied by Niederreiter \cite{rNIE95b,rNIE95c}.
The general recurrence has the form
\eq
  \bx_n = (A_1 \bx_{n-1} + \cdots + A_k \bx_{n-k}) \mod m  \label {MMRG}
\endeq
where $k$ and $m$ are the order and the modulus as for the MRG,
$\bx_n = (x_{n,1},\dots,x_{n,w})'$ is a $w$-dimensional vector,
and each $A_j$ is a $w\times w$ square matrix, for some positive 
integer $w$.  The case $w=1$ corresponds to the usual MRG.
The recurrence (\ref{MMRG}) has full period $m^{kw}-1$ if and only if
$m$ is prime and the characteristic polynomial
$$
  f(x) = \det\left( x^k I - x^{k-1} A_1 - x^{k-2} A_2 - \cdots - A_k\right)
$$
is a primitive polynomial modulo $m$ \cite{rNIE95b}.

There are different ways of producing the output.
We consider the following 3 cases:
\begin {eqnarray}
 u_{nw+i} &=& x_{n,i} / m \qquad\mbox{ for $0\le i<w$ and $n\ge 0$},
                                                      \eqlabel {MMRG:out1} \\
 u_{n}    &=& {1\over m} \left(\sum_{i=1}^w b_i x_{n,i} \mod m\right) 
           = \sum_{i=1}^w b_i x_{n,i} / m \mod 1
                        \qquad\mbox{ for $n\ge 0$},   \eqlabel {MMRG:out2} \\
 u_{n}    &=& \sum_{i=1}^w x_{n,i} m^{-i} 
                        \qquad\mbox{ for $n\ge 0$},   \eqlabel {MMRG:out3}
\end {eqnarray}
where $b_1,\dots,b_w$ are positive integers.
Case 3 is that used in \cite{rNIE95b} and does not give
rise to a lattice structure for $\Psi_t$ in the usual sense.
For both cases 1 and 2, the set $\Psi_t$ is the intersection of a lattice
$L_t$ with the unit hypercube.
Case 1 is used in \cite{rNIE95c}, where pseudorandom numbers are
generated in vector form, $w$ at a time.
It is also explained in \cite{rNIE95c} how to construct a basis for
the lattice $L_t$ when $t$ is a multiple of $w$.
(The generalization to other values of $t$ is trivial.)

%  --- Give equivalence with MRG.


%%%%%%%%%%%%%%%%%%%%%%%%%%%%%%%%%%%%%%%%%%%%%
\subsection {Multiply-with-carry}

A {\em Multiply-with-Carry\/} (MWC) generator 
\cite{rCOU95a,rCOU97a,rKOC95a,rMAR94a} is based on the recurrence
\begin{eqnarray}
  x_n &=& (a_1 x_{n-1} + \cdots + a_k x_{n-k} + c_{n-1}) \mod b,
                                                          \eqlabel {mwc1} \\
  c_n &=& (a_1 x_{n-1} + \cdots + a_k x_{n-k} + c_{n-1}) \div b,
                                                          \eqlabel {mwc2} \\
  u_n &=& x_n / b. \nonumber
\end {eqnarray}
where ``div'' denotes the integer division.
The recurrence looks like that of an MRG, except
that a {\em carry\/} $c_n$ is propagated between the steps.

Assume that $b$ is a power of 2, define $a_0 = -1$,
 $$ m = \sum_{\ell=0}^k a_\ell b^\ell, $$
and let $a$ be the inverse of $b$ in arithmetic modulo $m$.
For simplicity, assume $m > 0$.
Then, up to precision $1/b$, the MWC generator is equivalent to the LCG:
\eq
  z_n = a z_{n-1} \mod m;  \qquad  w_n = z_n/m.             \eqlabel {lcgz}
\endeq
In other words, if 
\eq
  w_n = \sum_{i=1}^\infty x_{n+i-1} b^{-i}                \eqlabel {wdigit}
\endeq
holds for $n=0$, then it holds for all $n$, and consequently
$|u_n - w_n| \le 1/b$ for all $n$.
The (approximate) lattice structure of the MWC can therefore be 
analyzed by analyzing that of the corresponding LCG (\ref{lcgz}).
This is what the LatMRG package does.

If $a_\ell \ge 0$ for $\ell\ge 1$, then all the recurrent states 
of the MWC satisfy $0\le c_n < a_1 + \cdots + a_k$.
In view of this inequality, we want the $a_\ell$ to be small,
so that their sum fits into a computer word
(e.g., $a_1 + \cdots + a_k \le b$).
But the coefficients should not be too small either, because
in dimension $t = k+1$, one has (see \cite{rCOU97a}):
\eq
  \ell_t = (1 + a_1^2 + \cdots + a_k^2)^{1/2}.            \eqlabel {dtkp1}
\endeq

Since $b$ is a power of 2, $a$ is a quadratic residue and so
cannot be primitive mod $m$.  Therefore the period length cannot
reach $m-1$ even if $m$ is prime.
But if $(m-1)/2$ is odd and 2 is primitive mod $m$
(e.g., if $(m-1)/2$ is prime), then (\ref{lcgz}) has
period length $\rho = (m-1)/2$.

%%%%%%%%%%%%%%%%%%%%%%%%%%%%%%%%%%%%%%%%%%%%%
\subsection {Combined generators}

\iffalse
Combining LCGs or MRGs with distinct moduli provides a efficient way
of implementing linear recurrences based on larger (non-prime) moduli.
The combined generator is equivalent, or can be approximated (depending
on the combination approach), by an associated LCG or MRG 
\cite{rLEC91b,rLEC96b}.
The package {\tt LatMRG} permits one to specify combined generators
in terms of their components.
Since these combined generators have composite moduli, they cannot
have full period.  In many cases, they also possess transient states.
In the latter case, one can decide to analyze either the lattice $L_t$
generated by all possible initial states, or that generated by the set
of recurrent states (see {\tt LatticeType} in {\tt MRG} and in {\tt seekl}).
\fi

Combining LCGs or MRGs with relatively prime moduli provides a efficient way
of implementing linear recurrences based on larger (non-prime) moduli.
The combination method that we consider adds, modulo $1$, the outputs of the 
components.
The package {\tt LatMRG} permits one to specify a {\em product} 
MRG in terms of component MRGs with relatively prime moduli. 
Its modulus is the product of the component moduli and its order is 
the maximum of the orders of the components. The recurrence governing
this product MRG, when taken modulo any one of the component moduli,
reduces to the component recurrence.
The combined generator can then be studied via this product generator, 
since one can view the former as embedded in the latter, and since
both have the same set of recurrent states (see \cite{rCOU96a}).
Facilities are provided to analyze, for any given MRG, 
either the lattice $L_t$
generated by all possible initial states, or that generated by the set
of recurrent states (see {\tt LatticeType} in {\tt MRG} and in {\tt seekl}).

Other types of combinations that have been proposed in the literature
are (often, depending on the parameters) closely approximated by
combinations of the above types \cite{rLEC91b,rLEC96b}.
They can thus be analyzed with the present software.

%%%%%%%%%%%%%%%%%%%%%%%%%%%%%%%%%%%%%%%%%%%%%%
\subsection {Computing a shortest nonzero vector or a reduced basis}

The module {\tt REDBAS} computes a shortest nonzero vector in a
lattice via the branch-and-bound (BB) algorithm proposed by \cite{mFIN85a},
with some additional refinements. 
For large dimensions $t$, this algorithm is much faster than the
algorithm given in \cite {rDIE75a,rKNU81a}.
The module also computes a MRLB via the algorithm of \cite{rAFF85a},
which works by successive applications of the BB procedure 
for finding a shortest vector.
The bounds in the BB procedure are computed through
a Choleski decomposition performed in (double precision)
floating-point arithmetic.
Numerical roundoff errors occur during these computations
and could (eventually) affect the results:
Because of slightly wrong bounds in the BB, one may miss a shorter
vector and, as a result, (conceivably) not obtain a true MRLB at the 
end of the reduction algorithm.
The module {\tt REDBAS2} perform basis reductions
by taking into account all sources of numerical error during
the computations and computing error bounds,
thereby yielding ``guaranteed error-free'' results.
Of course, this implies a lot of overhead, which means that the
procedures of {\tt REDBAS2} are much slower than
those of {\tt REDBAS}.
Another approach is used in module {\tt REDBAS3} bypassing all floating
point calculations. Again this is slower than {\tt REDBAS}.

%%%%%%%%%%%%%%%%%%%%%%%%%%%%%%%%%%%%%%%%%%%%%%%%%%%%%%%%%
\subsection {Representation of large numbers}

{\tt LatMRG} can deal with very large moduli and multipliers.
There is no limit on size other than the size of the computer memory
(and the CPU time).
For example, a generator with a modulus of a few hundred bits can be
analyzed easily.
Operations on large integers are performed using the package SENTIERS
\cite {iLEC88d}, also written in Modula-2.
Of course, these operations are performed in software and are significantly
slower than the standard operations supported by hardware.
For this reason, most of the basic (low-level) operations required by our
higher-level modules have been implemented in two versions.

When building a basis or checking maximal period conditions,
the modulus and multipliers can be represented either as {\tt LONGINT}
(32-bit integers) or {\tt SuperInteger} (arbitrary large integers, from
SENTIERS). The modules {\tt MULTLI} and {\tt MULTSI} support these
two representations, and provide basic facilities to the other modules
for manipulating the multipliers.
These are two versions of the generic module {\tt MULT}.

After a lattice basis and its dual have been constructed, when working 
on the basis (finding a shortest vector, Minkowski reduction, etc.),
the vector elements can be represented either as {\tt LONGREAL}
(64-bit floating-point numbers) or {\tt SuperInteger}.
The modules {\tt BASISLR} and {\tt BASISSI} implement basic facilities 
for these two cases. They are two versions of the generic module {\tt BASIS}.

In summary, {\tt MULT} is for working on the construction of a lattice basis
and its dual from the specifications of the generator, lacunary indices, 
etc., whereas {\tt BASIS} is for working on the basis (reduction, etc.)
after it has been constructed.

For the computations that are performed in floating-point arithmetic,
one can compute error bounds on everything (using {\tt REDBAS2}),
or no error bounds (using {\tt REDBAS}).
When performing a search for good generators, for instance, one can
first perform all the ``screening'' computations (involving many generators)
without computing the error bounds, and then recompute (verify)
with the error bounds only for the retained generator(s).

%%%%%%%%%%%%%%%%%%%%%%%%%%%%%%%%%%%%%%%%%%%
\subsection {Current implementations}

The software is currently working under the XDS Modula-2 environment,
on SUN computers under Solaris and on PC-type computers under Linux.
%% the MCS Modula-2 system from ModulaWare \cite{iJAG92a}.
%% with the SUN Modula-2 compiler.
%%  under VAX/VMS, with the compiler
%%  described in \cite{sECK85a}, under MS-DOS in the {\em TopSpeed}
%%  environment \cite {sKIN88a}, and
It can be ported to other machines for which a Modula-2 compiler is available.
Its implementation is built upon two other Modula-2 libraries called
{\tt mylib} \cite{iLEC01m} and SENTIERS \cite{iLEC88d}, 
which must also be ported for {\tt LatMRG} to work on other platforms.





